
\chapter{Part VI Review: Scales and Locality}

Renormalization is the statement that short-distance sensitivity can be absorbed into local terms already present, or into higher-dimension terms organized by an effective expansion.

\section{A guided reconstruction}

% QFTFIGURE-BEGIN partreview06-01-part-map
\qftfigure{figures/instructional/partreview06/partreview06-01-part-map.pdf}{The map collects the part's main constructions into a visual route so the reader can see how the chapters support one another.}{fig:partreview06-01-part-map}
% QFTFIGURE-END partreview06-01-part-map












A divergent loop integral is not yet an observable.  A measured condition fixes a renormalized parameter, and a counterterm absorbs regulator dependence.  If the subtraction scale changes while the observable is held fixed, the coupling runs:
\[
  \mu\frac{\dd g}{\dd\mu}=\beta(g).
\]
Wilsonian language makes the same idea physical by integrating out high-momentum modes and watching the coefficients of operators change.

\begin{exercise}
Explain the difference between a regulator, a counterterm, and a renormalization condition.
\begin{answer}
A regulator makes an integral defined, a counterterm absorbs regulator dependence, and a renormalization condition ties the finite parameter to a measured quantity.
\end{answer}
\end{exercise}
